% Compile from this directory:
%   pdflatex main
%   bibtex main
%   pdflatex main
%   pdflatex main

\documentclass[11pt]{article}

\usepackage{amsmath,amssymb,amsthm,mathtools}
\usepackage{geometry}
\geometry{margin=1in}
\usepackage{microtype}
\usepackage{graphicx}
\usepackage{xcolor}
\usepackage[colorlinks=true,linkcolor=blue!65!black,citecolor=green!45!black,urlcolor=blue!65!black,
  pdftitle={A matrix-valued von Mangoldt measure in the finite Connes-van Suijlekom path},
  pdfauthor={Akiva Groskin},
  pdfsubject={Riemann zeta function; Weil explicit formula; finite Weil form; von Mangoldt measure},
  pdfkeywords={Riemann zeta function, Weil explicit formula, von Mangoldt function, Connes-van Suijlekom, finite matrix path}]{hyperref}
\usepackage{url}
\urlstyle{tt}
\emergencystretch=1.5em

\theoremstyle{plain}
\newtheorem{theorem}{Theorem}[section]
\newtheorem{proposition}[theorem]{Proposition}
\newtheorem{lemma}[theorem]{Lemma}
\newtheorem{corollary}[theorem]{Corollary}

\theoremstyle{definition}
\newtheorem{definition}[theorem]{Definition}

\theoremstyle{remark}
\newtheorem{remark}[theorem]{Remark}

\newcommand{\R}{\mathbb R}
\newcommand{\Z}{\mathbb Z}
\newcommand{\cvs}{Connes--van Suijlekom}
\newcommand{\ccm}{Connes--Consani--Moscovici}
\newcommand{\dd}{\,d}
\newcommand{\one}{\mathbf 1}
\newcommand{\diag}{\operatorname{diag}}
\newcommand{\adj}{\operatorname{adj}}
\newcommand{\sgn}{\operatorname{sgn}}
\newcommand{\Lam}{\Lambda}

\title{A matrix-valued von Mangoldt measure\\
in the finite \cvs\ path}

\author{Akiva Groskin}

\date{July 2026}

\begin{document}

\maketitle

\begin{abstract}
Fix a Galerkin level $N$ in the finite \cvs\ truncation of the Weil
quadratic form (with no archimedean cutoff), and let $u=\log c$ be the prime
cutoff, which is retained and varied. The object of this paper is the singular
part of $D_u^2Q_N$ along the finite matrix path $u\mapsto Q_N(u)$: a negative
semidefinite rank-one \emph{matrix-valued von Mangoldt measure}, a finite,
cutoff-free realization of the prime side of the Weil--Guinand explicit
formula on the operator path itself. Its defining computation is the
first-derivative jump at a prime-power threshold $u=\log q$, which returns the
von Mangoldt weight exactly, in every matrix entry:
\[
        Q_N'(\log q+0)-Q_N'(\log q-0)
        =-\frac{2\Lam(q)}{\sqrt q\,\log q}\,\one_N\one_N^t .
\]
This identity is an elementary structural derivative of the path's defining
prime sum: the jump is one differentiation of the entering prime atom, and
that is precisely what makes the readout exact rather than asymptotic; the
identity relating the prime side to the zeros is classical and is not proven
here. The contribution is the isolation and naming of the finite matrix-valued
measure and the exact finite geometry established around it. First, the event
is \emph{arithmetically rigid}:
the entire singular jet at an edge is a universal matrix scaled by the single
scalar $\Lam(q)$, so it carries exactly the von Mangoldt weight and nothing
more. Second, the event geometry is a sharp piece of finite harmonic analysis: the
source-to-jet map is a lossless finite dictionary on the squared-integer nodes
$\{0,1^2,\dots,N^2\}$, with an explicit confluent-Vandermonde determinant, a
sharp $2N{+}1$ window, and a universal recurrence $S\prod_{k=1}^N(S-k^2)^2$; a
sharp finite vanishing-moment bound (a \emph{prime-edge uncertainty principle}
in the band-limited sense) bounds the order at which a
band-limited normal can be blind to an entering prime, with the centered
finite-difference stencil the unique extremizer; and each prime event is a
rank-one negative-semidefinite deflation of the spectral velocity, interlacing
the velocity spectrum. The statements are finite-dimensional and proved under
explicitly stated hypotheses; the event algebra transfers to any $L$-function of
the Selberg class once a finite path with the same event structure is prescribed
(the arithmetic entering only through the source weights), while the
sign-definite consequences require a positive weight and so hold for $\zeta$
(the spectral-barrier statement remaining conditional on its
positive-definiteness hypothesis). The results are structural; the paper
proves no positivity, no Riemann Hypothesis, no prime-counting, next-prime, or
factoring statement.
\end{abstract}

\medskip
\noindent\textit{2020 Mathematics Subject Classification.} Primary 11M26;
Secondary 11M06, 15A15, 47A75, 42C40.

\noindent\textit{Key words and phrases.} Riemann zeta function, Weil explicit
formula, von Mangoldt function, Connes--van Suijlekom operator, truncated Weil
form, confluent Vandermonde, vanishing moments.

\medskip
\noindent\textit{Note on this version.} This version corrects version 1 of the
Zenodo deposit: the coincidence-averaging corollary now carries its
uncorrelatedness hypothesis explicitly (with the covariance discussion in
Remark~\ref{rem:covariance}); the directional-residual and singular-ratio
diagnostics are no longer called equivalent; the coincidence-resolvent pole
statement is stated by spectral projection, allowing removable points; the
reciprocal Weyl identity is stated meromorphically with its pointwise domain
made exact; and the Krein-string boundary-mass reading is downgraded from a
theorem to an explicitly labeled analogy. Scope language is also recalibrated
relative to version 1: the abstract's exactness sentence is removed, the
Selberg-class transfer is conditioned on an actually constructed finite path,
the spectral-barrier statement is demoted to a conditional remark stating its
positive-definiteness hypothesis, one proof argument in the resolution-scale
corollary is repaired (a domination argument replaces a first-positive-root
phrase), the verification table marks the variance reduction as an IID-model
statement, the vanishing-moment result is presented under a calibrated name
(with the uncertainty-principle reading kept as a gloss), the related-work
discussion is extended (Makraini, Andrade, Andrews), and point-of-use
attributions are added. The central rank-one prime-power
jump, the distributional matrix-valued measure, the dictionary,
vanishing-moment, and deflation results are unchanged in statement.

\section{Introduction}

The \cvs\ and \ccm\ constructions attach finite Galerkin matrices to the
truncated Weil quadratic form
\cite{ConnesvanSuijlekom2025,ConnesConsaniMoscovici2025}. Paper~1 of this
sequence implemented the finite \cvs\ computation and zero recovery numerically
\cite{Groskin2026Paper1}; Paper~2 isolates the cutoff-free finite dictionary
from Galerkin coefficient vectors to compactly supported Guinand--Weil test
functions and proves the finite archimedean tail-order rule
\cite{Groskin2026FiniteDictionary}. That cutoff-free finite matrix is the object
studied here. We fix $N$ and vary the prime cutoff $c=e^{u}$.

The most concrete fact is the following. At a prime-power threshold $u=\log q$,
$q=p^{a}$, a new prime source enters the finite prime sum. It is invisible in
the matrix value itself, but it is exactly visible in the first derivative, and
\emph{any one matrix entry recovers the von Mangoldt weight},
\[
        \Lam(q)=-\frac{\sqrt q\,\log q}{2}\,
        \bigl(Q_N'(\log q+0)-Q_N'(\log q-0)\bigr)_{mn}
        \qquad\text{for every }m,n.
\]
This identity is elementary and holds by construction: the prime sum is a
defining summand of $Q_N$, and the jump is one chain-rule differentiation of
the entering atom (Section~\ref{sec:measure}). The exactness is the point, and
the contribution of the paper is what is isolated around it: the finite
matrix-valued measure itself, its rigidity, the exact squared-node dictionary,
and the sharp vanishing-moment ceiling with its unique extremizer.
Summing these jumps against a test function reproduces $\sum_{q}\Lam(q)\phi(\log
q)$, the prime side of the Weil--Guinand explicit formula
\cite{Weil1952,IwaniecKowalski2004}; the archimedean and pole terms
(Paper~2) carry the complementary side, and the eigenvalues of $Q_N$
approximate the zeros (Paper~1; a quantitative convergence law for the first
zero in the CCM setting is given by Andrews \cite{Andrews2026}). In this sense the finite path renders the
\emph{prime side} of the explicit formula exactly and without an archimedean
cutoff, with the primes appearing as the singular part of a second derivative;
the archimedean and pole side is closed-form (Paper~2) and the zeros enter only
approximately (Paper~1). Morán Ledezma gives a
probabilistic ``arithmetic spectral measure'' reading of the explicit sums
\cite{MoranLedezma2023}; the object here is the finite matrix-valued measure
carried by the \cvs\ path, distinct from that scalar construction.

The rest of the arithmetic story is a rigidity theorem: beyond the scalar
$\Lam(q)$, the local event contains no further information about $q$, and the
universal singular jet has an explicit closed form (Section
\ref{sec:rigidity}). The redundancy of the rank-one measure has a concrete
payoff: under mean-zero, equal-variance, pairwise uncorrelated entrywise
errors, averaging its $(2N+1)^2$
entries reads $\Lam(q)$ with variance reduced by
the same factor (Corollary~\ref{cor:matched}; the covariance hypothesis is
essential, Remark~\ref{rem:covariance}). What remains is the finite algebra of
the event geometry, recorded as a reference catalog. The
finite source-to-jet map is a lossless dictionary whose nodes are the squares
$0,1^{2},\dots,N^{2}$ (Section~\ref{sec:dictionary}); a sharp vanishing-moment
ceiling governs how blind a band-limited probe can be to an entering prime,
with a unique extremizer (Section~\ref{sec:uncertainty}); and the rank-one
event deflates the spectral velocity, interlaces its spectrum, and is generated
by a single coincidence resolvent (Section~\ref{sec:deflation}). The event
algebra transfers across the Selberg class conditionally on the corresponding
finite paths, with a residue-class readout of
$\Lam(q)$ across the Dirichlet family under the same hypothesis; the
sign-definite parts require a positive
weight (Section~\ref{sec:scope}).

Here an edge is treated only after a nonsmooth point of the path has been
localized; nothing below predicts future prime powers from sampled data, and
nothing below is a positivity statement. To our knowledge the finite \cvs\
matrix-valued von Mangoldt measure and its rigidity, together with the
squared-node event dictionary and the finite prime-edge vanishing-moment
ceiling,
have not previously been isolated. The trace-formula setting is due to Connes
\cite{Connes1999,Connes2026}; the finite Galerkin framework to \cvs\ and \ccm;
Suzuki develops a continuous screw-function framework
\cite{Suzuki2026}; Śliwiński studies the adjacent $D_{\log}^{(\lambda,N)}$
spectral-error regime \cite{Sliwinski2026}; Makraini records unconditional
spectral bounds, a Krein-space symmetry, and a norm-convergence obstruction
for the truncated Weil operator compressed to a prolate Sonin space, reducing
positivity in that framework to an asymptotic kernel-invisibility property
\cite{Makraini2026}, a line concerning positivity and convergence in the
cutoff rather than the derivative structure of a fixed-$N$ path; Andrade records related Loewner,
exact-entry, operator-monotone, and rank-one pole-resolvent framings
\cite{Silva20710075,Silva20737111,Silva20671635,Silva20682834}; and Liao uses prime-event
language for statistical event sequences \cite{Liao2026}. The classical Prony
and confluent-Vandermonde algebra
\cite{Prony1795,Gautschi1962,BatenkovYomdin2013}, the explicit formula
\cite{Weil1952,IwaniecKowalski2004}, and the vanishing-moment (sum-rule)
inequality \cite{Daubechies1992} are used as tools; what is new, to our
knowledge, is their exact
realization on the finite \cvs\ path.

\section{The finite path}\label{sec:path}

Fix $N\ge0$ and write
\[
        I_N=\{-N,\ldots,N\},\qquad d=2N+1,\qquad
        \one_N=(1)_{m\in I_N}.
\]
For a single prime source at reduced position $\omega\in[0,1]$, the finite
Galerkin matrix of that source is the divided-difference matrix $A_N(\omega)$ on
$I_N$,
\[
 (A_N(\omega))_{mn}=
 \begin{cases}
 \dfrac{\sin(2\pi m\omega)-\sin(2\pi n\omega)}{\pi(m-n)},&m\ne n,\\[1.2ex]
 2\omega\cos(2\pi m\omega),&m=n,
 \end{cases}
\]
in the normalization of the finite dictionary \cite{Groskin2026FiniteDictionary}.
A symbol-level Loewner divided-difference representation of the prime
contribution in the localized Weil form was recorded earlier by Andrade
\cite{Silva20710075}; what is used here is the derivative jump of the path
across the threshold, not the representation itself.
Its off-diagonal entries are divided differences of $m\mapsto\sin(2\pi m\omega)$;
this is why a source entering at $\omega=0$ is invisible in the matrix value but
exact in the first derivative (Lemma~\ref{lem:edge}).

\begin{remark}[Divided-difference structure]\label{rem:dd}
Write $\mathrm{DD}_X[g]$ for the confluent divided-difference matrix of a $C^1$
function $g$ on real nodes $X=\{x_m\}$, with entries
$(g(x_m)-g(x_n))/(x_m-x_n)$ off the diagonal and $g'(x_m)$ on it. Then
$A_N(\omega)=\mathrm{DD}_X[g_\omega]$ with $X=\{-N,\dots,N\}$ and
$g_\omega(x)=\sin(2\pi x\omega)/\pi$: the diagonal $2\omega\cos(2\pi m\omega)$ is
exactly the confluent limit $g_\omega'(m)$. This is the divided-difference
representation of the finite quadratic form established in
\cite[Props.~4.1--4.2]{ConnesvanSuijlekom2025}; here it is specialized to a
single prime source and differentiated in the prime cutoff along the path. The
whole event package below uses
only this structure. Indeed, for any real-analytic background $R(u)$, amplitude
$s(u)$, and reduced position $\omega(u)$ with $\omega(u_0)=0$,
$\omega'(u_0)\ne0$, such that $g_0$ is constant on the node span (so
$\mathrm{DD}_X[g_0]=0$) and $\tfrac{d}{d\omega}g_\omega|_0$ is affine on the nodes with
nonzero slope $\lambda$, the path
$u\mapsto R(u)-s(u)\,\mathrm{DD}_X[g_{\omega(u)}]\,\mathbf 1[u\ge u_0]$, with the
source term present only for $u\ge u_0$ (the entering-source mechanism of
Theorem~\ref{thm:jump}), has the
universal rank-one first-derivative jump $-s(u_0)\omega'(u_0)\lambda\,
\one_N\one_N^t$ and the ensuing rigidity, independent of the node geometry. The
finite \cvs\ path is the case $g_\omega=\sin(2\pi x\omega)/\pi$ ($g_0=0$,
$\lambda=2$),
$s=\Lam(q)/\sqrt q$, $\omega=1-\log q/u$.
\end{remark}

Let $Q_N(u)$ denote the cutoff-free finite \cvs/\ccm\ matrix at prime cutoff
$c=e^u$:
\[
        Q_N(u)=Q_{\mathrm{arch}}(u)+Q_{\mathrm{pole}}(u)
        -\sum_{q=p^a\le e^u}\frac{\Lam(q)}{\sqrt q}\,
             A_N\!\left(1-\frac{\log q}{u}\right).
\]
Here $Q_{\mathrm{arch}}$ and $Q_{\mathrm{pole}}$ are the closed-form finite
archimedean and pole matrices of the finite dictionary, given explicitly in
Paper~2 \cite[Lemma~2.1 and the displayed pole and archimedean
sources]{Groskin2026FiniteDictionary}:
$Q_{\mathrm{pole}}$ is the divided-difference matrix of the pole source
$\psi_0$, and $Q_{\mathrm{arch}}$ that of the archimedean density, both in the
same $A_N$-normalization as the prime blocks. We call the path
\emph{archimedean-cutoff-free} because no archimedean cutoff is present: only the
prime cutoff $e^u$ is varied. This is the standard finite \cvs/\ccm\ assembly
with the prime cutoff exposed as a variable, not an ad hoc signal. From Paper~2
we use only the following regularity of the non-prime blocks.

\begin{lemma}[Background regularity, imported from Paper~2]\label{lem:background}
The archimedean and pole blocks $Q_{\mathrm{arch}}(u)$ and $Q_{\mathrm{pole}}(u)$
are real analytic on $(0,\infty)$. Consequently the only points at which
$u\mapsto Q_N(u)$ can fail to be real analytic are the prime-power thresholds
$u=\log q$.
\end{lemma}

\begin{proof}
This lemma is imported: the blocks are defined, entrywise and in closed form,
in \cite[Lemma~2.1]{Groskin2026FiniteDictionary}, and the analyticity below
only restates the regularity of those closed forms.
Entrywise, the pole block is a rational-hyperbolic expression whose denominators
$u^2+16\pi^2n^2$ are strictly positive for $u>0$, and the archimedean block is
built from digamma and trigamma values at $\tfrac14+\tfrac{i\pi n}{u}$, whose
arguments avoid the poles of $\psi,\psi'$ for $u>0$, together with a correction
series of geometric ratio $e^{-2u}$. Each summand extends holomorphically to a
complex neighborhood of $(0,\infty)$, and on such a neighborhood the $e^{-2u}$
ratio bound gives uniform convergence on compact subsets, so each block is real
analytic there \cite{Groskin2026FiniteDictionary}. Away from the finitely many
thresholds in a compact interval, each reduced position $1-\log q/u$ is real
analytic in $u$, so $Q_N$ is real analytic off $\{\log q\}$.
\end{proof}

\begin{lemma}[Edge derivative]\label{lem:edge}
The single-source matrix satisfies $A_N(0)=0$ and $A_N'(0)=2\,\one_N\one_N^t$.
\end{lemma}

\begin{proof}
The diagonal gives $A_N(0)_{mm}=0$, and for $m\ne n$ the numerator vanishes at
$\omega=0$, so $A_N(0)=0$. Differentiating at zero gives
$(2\pi m-2\pi n)/(\pi(m-n))=2$ for $m\ne n$, and the diagonal derivative of
$2\omega\cos(2\pi m\omega)$ at zero is $2$.
\end{proof}

A companion structural fact, used throughout Section~\ref{sec:dictionary}, is
that every source profile has only odd Taylor content at the edge.

\begin{lemma}[Odd content]\label{lem:odd}
For each $m,n$ the entry $\omega\mapsto(A_N(\omega))_{mn}$ is an odd function of
$\omega$; its even Taylor coefficients at $\omega=0$ vanish, and all arithmetic
content of the event sits in the odd jets.
\end{lemma}

\begin{proof}
Both $\sin(2\pi m\omega)$ and $\omega\cos(2\pi m\omega)$ are odd in $\omega$, and
the off-diagonal quotient of odd functions by the constant $\pi(m-n)$ is odd.
\end{proof}

\section{The matrix-valued von Mangoldt measure}\label{sec:measure}

\begin{theorem}[Prime-power derivative jump]\label{thm:jump}
The path $u\mapsto Q_N(u)$ is continuous and piecewise real analytic on
$(0,\infty)$, with first-derivative jumps only at $u=\log q$, $q=p^a$. At such a
point,
\[
        \Delta Q_N'(\log q):=Q_N'(\log q+0)-Q_N'(\log q-0)
        =-\frac{2\Lam(q)}{\sqrt q\,\log q}\,\one_N\one_N^t,
\]
and consequently, for every entry $(m,n)$,
\[
        \Lam(q)=-\frac{\sqrt q\,\log q}{2}
        \bigl(Q_N'(\log q+0)-Q_N'(\log q-0)\bigr)_{mn}.
\]
\end{theorem}

\begin{proof}
By Lemma~\ref{lem:background} the path is real analytic off $\{\log q\}$, hence
piecewise real analytic. At $u_0=\log q$ the entering source has
$\omega_q(u)=1-\log q/u$, so $\omega_q(u_0)=0$ and $\omega_q'(u_0)=1/\log q$.
Since $A_N(0)=0$ the path is continuous at $u_0$, and the derivative jump comes
only from the entering atom and Lemma~\ref{lem:edge}:
\[
 -\frac{\Lam(q)}{\sqrt q}\,A_N'(0)\,\omega_q'(u_0)
 =-\frac{2\Lam(q)}{\sqrt q\,\log q}\,\one_N\one_N^t .
\]
All other terms are analytic at $u_0$ and cancel. The recovery formula follows
because every entry of $\one_N\one_N^t$ equals $1$.
\end{proof}

\noindent Theorem~\ref{thm:jump} is an elementary structural derivative of the
path's defining prime sum: the prime block is a defining summand of $Q_N(u)$,
and the jump is one chain-rule differentiation of the entering atom
(Lemma~\ref{lem:edge}). This is a feature, not an accident: it is what makes
the readout exact rather than asymptotic. The contribution of the paper is not
the differentiation but the isolation of the resulting measure
(Corollary~\ref{cor:measure}) and the exact finite geometry around it
(Sections~\ref{sec:rigidity}--\ref{sec:deflation}); see the scope discussion in
Section~\ref{sec:scope}.

\begin{corollary}[Matrix-valued von Mangoldt measure]\label{cor:measure}
In the sense of distributions on $(0,\infty)$,
\[
        (D_u^2Q_N)_{\rm sing}
        =-\sum_{q=p^a}\frac{2\Lam(q)}{\sqrt q\,\log q}\,
        \one_N\one_N^t\,\delta_{\log q}.
\]
Each atom is negative semidefinite and has rank one. Writing $\mu_{mn}$ for the
scalar singular measure $\bigl((D_u^2Q_N)_{mn}\bigr)_{\rm sing}
=-\sum_q\frac{2\Lam(q)}{\sqrt q\,\log q}\,\delta_{\log q}$ (the same for every
$(m,n)$), one has, for every compactly supported $\phi$,
\[
 \sum_{q=p^a}\Lam(q)\phi(\log q)
 =-\frac12\int_0^\infty e^{u/2}u\,\phi(u)\,\dd\mu_{mn}(u).
\]
\end{corollary}

\begin{proof}
A continuous piecewise-analytic scalar whose first derivative jumps by $J$ at
$u_0$ has singular second derivative $J\delta_{u_0}$; apply this entrywise to
Theorem~\ref{thm:jump}. At $u=\log q$ the factor $-e^{u/2}u/2$ converts
$2\Lam(q)/(\sqrt q\log q)$ back to $\Lam(q)$.
\end{proof}

Corollary~\ref{cor:measure} is the paper's central object: the finite \cvs\ path
carries the prime side of the explicit formula \cite{Weil1952,IwaniecKowalski2004}
as an exact, cutoff-free, negative semidefinite rank-one \emph{matrix-valued von
Mangoldt measure}. Figure~\ref{fig:event} shows a scalar readout of the path and
the corresponding measure of rank-one shocks.

\begin{figure}[t]
\centering
\includegraphics[width=\textwidth]{fig_event_signal.pdf}
\caption{The finite path as an arithmetic event signal (top: a scalar entry of
the prime part, continuous and piecewise real analytic, with a slope kink at
each $u=\log q$) and its second-derivative singular part (bottom: the
matrix-valued von Mangoldt measure, a rank-one negative shock of height
$2\Lam(q)/(\sqrt q\,\log q)$ at each prime power, including the true powers
$4=2^2$, $8=2^3$, $9=3^2$).}
\label{fig:event}
\end{figure}

The first derivative jumps down by a rank-one shock; the second derivative jumps
\emph{up} by the complementary rank-one shock, in the same universal direction.

\begin{proposition}[Second-order event law]\label{prop:second}
At a prime-power event $u_0=\log q$ the classical (non-singular) second derivative
of the path has the rank-one positive semidefinite jump
\[
        Q_N''(\log q+0)-Q_N''(\log q-0)
        =+\frac{4\Lam(q)}{\sqrt q\,(\log q)^2}\,\one_N\one_N^t,
\]
opposite in sign to the first-derivative jump of Theorem~\ref{thm:jump}.
\end{proposition}

\begin{proof}
By Lemma~\ref{lem:background} the jump comes only from the entering atom
$E_q(\tau)=-\tfrac{\Lam(q)}{\sqrt q}A_N(\omega_q(\tau))$, $\omega_q=\tau/(u_0+\tau)$.
Then $E_q''(0+)=-\tfrac{\Lam(q)}{\sqrt q}\bigl(A_N''(0)\,\omega_q'(u_0)^2+
A_N'(0)\,\omega_q''(u_0)\bigr)$; since $A_N$ is odd (Lemma~\ref{lem:odd})
$A_N''(0)=0$, and with $A_N'(0)=2\one_N\one_N^t$, $\omega_q''(u_0)=-2/u_0^2$ this is
$+\tfrac{4\Lam(q)}{\sqrt q\,u_0^2}\one_N\one_N^t$.
\end{proof}

Because the jump lies entirely in the known rank-one direction $\one_N\one_N^t$,
every entry of the jump matrix carries the same scalar, so the $(2N+1)^2$
entries can be pooled. Since all entries of the model jump \emph{coincide} in
the universal direction, we call the resulting all-ones compressions
\emph{coincidence} objects: the coincidence average below, and the coincidence
resolvent and coincidence-Weyl function of Section~\ref{sec:deflation}. How
much pooling helps is decided by the correlation
structure of the measurement errors, which must therefore be part of the
hypothesis.

\begin{corollary}[Coincidence-averaged weight readout]\label{cor:matched}
Suppose the first-derivative jump is measured or estimated as
$\widehat{J}=-a_q\,\one_N\one_N^t+E$, $a_q=2\Lam(q)/(\sqrt q\,\log q)$, with
entrywise errors $E_{mn}$ of mean zero and common variance $\sigma^2$ that are
\emph{pairwise uncorrelated}. Among linear unbiased estimators of $a_q$, the
minimum-variance one (Gauss--Markov) is the coincidence average
\[
        \widehat a_q=-\frac{1}{(2N+1)^2}\sum_{m,n}\widehat{J}_{mn},
        \qquad \operatorname{Var}(\widehat a_q)=\frac{\sigma^2}{(2N+1)^2},
\]
a variance smaller by the factor $(2N+1)^2$ than the single-entry readout of
Theorem~\ref{thm:jump}. The residual $\widehat{J}+\widehat a_q\one_N\one_N^t$ is
orthogonal to $\one_N\one_N^t$, and its Frobenius size measures the distance of
$\widehat{J}$ to the one-dimensional model span of $\one_N\one_N^t$.
\end{corollary}

\begin{proof}
For the unit-Frobenius direction $u=\one_N\one_N^t/\|\one_N\one_N^t\|_F$, the
amplitude functional $\langle u,\cdot\rangle_F$ is the Gauss--Markov estimator of
the coefficient of $u$; with $\|\one_N\one_N^t\|_F^2=(2N+1)^2$ this is the stated
average, and for pairwise uncorrelated entries
$\operatorname{Var}=\sigma^2\|\one_N\one_N^t\|_F^2/(\|\one_N\one_N^t\|_F^2)^2
=\sigma^2/(2N+1)^2$, against $\sigma^2$ for one entry. The orthogonal split
$\widehat{J}=-\widehat a_q\one_N\one_N^t+R$, $R\perp\one_N\one_N^t$, gives the
residual statement.
\end{proof}

\noindent Corollary~\ref{cor:matched} is a design upper bound on the pooling
gain under an idealized error model, approximated for instance when the
entrywise estimates are formed from disjoint sample sets;
Remark~\ref{rem:covariance} records the realistic correlated regimes.

\begin{remark}[Covariance hypotheses and diagnostics]\label{rem:covariance}
(i) The uncorrelatedness hypothesis cannot be dropped: if all entries carry one
common error variable, $E_{mn}\equiv Z$ with $\operatorname{Var}(Z)=\sigma^2$,
then every mean/variance hypothesis holds while
$\operatorname{Var}(\widehat a_q)=\sigma^2$, with no reduction. For a general
\emph{nonsingular} (positive definite) error covariance $\Sigma$ the
minimum-variance linear unbiased readout is the
generalized least-squares estimator with weights $\Sigma^{-1}$, and its variance
is governed by $\Sigma$, not by the entry count; for singular $\Sigma$, as in
the common-error example just given, the best linear unbiased estimator exists
but is not given by this formula. (ii) If the measurement is
symmetric, $E_{mn}=E_{nm}$, the entries cannot all be uncorrelated; under the
natural model (the entries with $m\le n$, diagonal included, pairwise
uncorrelated, each of variance
$\sigma^2$) the coincidence average has variance $\sigma^2(2d-1)/d^{3}$ with
$d=2N+1$, a reduction of order $d^{2}/2$ rather than $d^{2}$: indeed
$\operatorname{Var}\bigl(\sum_{m,n}E_{mn}\bigr)
=d\,\sigma^2+4\binom d2\sigma^2=d(2d-1)\sigma^2$, since each of the $\binom d2$
off-diagonal values appears twice in the sum, and dividing by $d^4$ gives the
formula. (iii) The
directional residual and the singular-value ratio $\sigma_2/\sigma_1$ of
$\widehat J$ are \emph{different} diagnostics: the residual measures distance to
the fixed model direction $\one_N\one_N^t$, while $\sigma_2/\sigma_1$ measures
the relative (spectral-norm)
distance to the set of all rank-one matrices. A rank-one matrix in a different
direction has $\sigma_2/\sigma_1=0$ but a large directional residual, so
$\sigma_2/\sigma_1$ is only a necessary check; the model at the fixed direction
is certified by the directional residual.
\end{remark}

\noindent A one-sided difference estimate of the jump additionally carries the
systematic bias of Proposition~\ref{prop:second}; a centered difference removes it
to leading order.

\section{What the local event determines: arithmetic rigidity}\label{sec:rigidity}

The first jump recovers $\Lam(q)$. It is natural to ask whether the higher-order
local structure of the event carries any further arithmetic. It does not.

\begin{theorem}[Arithmetic rigidity of the event]\label{thm:rigidity}
Let $u_0=\log q$ be a prime-power edge. For every order $r\ge1$ the singular jet
of the path is
\[
        \Delta Q_N^{(r)}(u_0)=\frac{\Lam(q)}{\sqrt q}\,B_{r,N}(u_0),
\]
where $B_{r,N}(u_0)$ is a universal matrix depending only on $r$, $N$, and the
edge location $u_0$, and not otherwise on $q$. Consequently, modulo the known
edge location, the entire singular event content is the single scalar $\Lam(q)$:
at a fixed edge $u_0$ the assignment
$\Lam\mapsto(\Lam/\sqrt q)\,(B_{r,N}(u_0))_{r\ge1}$ is a linear injection, so
the singular jet family determines and is determined by the single scalar
$\Lam(q)$, and no
finer arithmetic invariant of $q$ is present in the local event.
\end{theorem}

\begin{proof}
By Theorem~\ref{thm:jump} and the analyticity of the background, all singular
jets at $u_0$ come from the entering atom
$E_q(\tau)=-\tfrac{\Lam(q)}{\sqrt q}A_N\!\bigl(\tau/(u_0+\tau)\bigr)$,
$\tau=u-u_0$: for each $r$,
\[
        \Delta Q_N^{(r)}(u_0)
        =-\frac{\Lam(q)}{\sqrt q}\,
        \frac{d^{\,r}}{d\tau^{\,r}}\Big|_{\tau=0}
        A_N\!\Bigl(\frac{\tau}{u_0+\tau}\Bigr)
        =\frac{\Lam(q)}{\sqrt q}\,B_{r,N}(u_0),
\]
and the bracket depends only on $r$, $N$, and $u_0$. Since $q=e^{u_0}$ is fixed
by the edge, $\sqrt q$ and every $B_{r,N}(u_0)$ are known; the jets are therefore
$\Lam(q)$ times a known matrix, and $\Lam(q)$ is read off from any one nonzero
entry of $B_{1,N}(u_0)=-2(\log q)^{-1}\one_N\one_N^t$.
\end{proof}

The argument is elementary, but the conclusion is structurally clean: it
delimits exactly what the local event can encode. The edge location $u_0=\log q$
already fixes $q$, and beyond it the singular jet adds only the scalar $\Lam(q)$;
the event is in this sense a rank-one arithmetic object. The remaining sections
show that the geometry \emph{around} the weight, though carrying no further
arithmetic, is a sharp and self-contained piece of finite harmonic analysis.

The universal matrices $B_{r,N}$ are not merely well defined; they have an
explicit closed form, which makes the rigidity statement concrete and reduces the
first two derivative laws to the cases $r=1,2$.

\begin{proposition}[Closed form of the universal jet]\label{prop:Bform}
Write $C_j=[\omega^j]A_N(\omega)$ for the $j$th matrix Taylor coefficient of the
single-source matrix; $C_j=0$ for even $j$, and for odd $j=2\ell+1$
\[
 (C_j)_{mn}=\frac{(-1)^\ell(2\pi)^j}{j!}\,\frac{m^j-n^j}{\pi(m-n)}\ \ (m\ne n),
 \qquad
 (C_j)_{mm}=\frac{2(-1)^\ell(2\pi m)^{\,j-1}}{(j-1)!}.
\]
Then the universal matrices of Theorem~\ref{thm:rigidity} are
\[
 B_{r,N}(u_0)=-\frac{r!}{u_0^{\,r}}
 \sum_{\substack{1\le j\le r\\ j\ \mathrm{odd}}}
 (-1)^{r-j}\binom{r-1}{j-1}\,C_j .
\]
In particular $B_{1,N}=-2u_0^{-1}\one_N\one_N^t$ (Theorem~\ref{thm:jump}) and
$B_{2,N}=+4u_0^{-2}\one_N\one_N^t$ (Proposition~\ref{prop:second}).
\end{proposition}

\begin{proof}
By Lemma~\ref{lem:odd}, $A_N(\omega)=\sum_{j\ \mathrm{odd}}C_j\,\omega^j$ with the
displayed coefficients, read off from the Taylor expansions of
$\sin(2\pi m\omega)/\pi$ and $2\omega\cos(2\pi m\omega)$. Compose with
$\omega(\tau)=\tau/(u_0+\tau)$ and use the edge transport
$[\tau^r]\omega^j=(-1)^{r-j}\binom{r-1}{j-1}u_0^{-r}$ (Lemma~\ref{lem:transport}
below); since $B_{r,N}=-r!\,[\tau^r]A_N(\omega(\tau))$ by
Theorem~\ref{thm:rigidity}, the sum follows. Evaluating at $r=1$ ($C_1=2\one_N\one_N^t$)
and $r=2$ recovers the first- and second-derivative jumps.
\end{proof}

\section{The finite source-to-jet dictionary}\label{sec:dictionary}

We record the finite algebra of the event geometry: how a general finite source
measure is encoded in the matrix, and how the edge jets invert that encoding.
Sections~\ref{sec:dictionary} and~\ref{sec:uncertainty} are a reference catalog
of this finite harmonic analysis, recorded for the event geometry itself.
The statements are exact facts about the fixed matrix $A_N$; by
Theorem~\ref{thm:rigidity} they specialize at a prime event to the recovery of a
fixed profile, not to further arithmetic.

Let $\mathcal V_N=\operatorname{span}\{b_0,b_k^s,b_k^c:1\le k\le N\}$ be the real
source-profile space, with
\[
 b_0(\omega)=\omega,\quad
 b_k^s(\omega)=\frac{\sin(2\pi k\omega)}{2\pi k},\quad
 b_k^c(\omega)=\omega\cos(2\pi k\omega).
\]
For $f=a b_0+\sum_k(A_kb_k^s+B_kb_k^c)$ write
$\mathfrak c_N(f)=(a,A_1,\dots,A_N,B_1,\dots,B_N)$.

\begin{lemma}[Finite source-matrix quotient]\label{lem:quotient}
For a finite signed measure $\mu$ on $[0,1]$ put $B_\mu=\int A_N\dd\mu$ and
\[
 M_0=\int\omega\dd\mu,\quad S_k=\int\sin(2\pi k\omega)\dd\mu,\quad
 C_k=\int\omega\cos(2\pi k\omega)\dd\mu.
\]
Then $\mu\mapsto B_\mu$ factors exactly through the $2N+1$ coordinates
$(M_0,S_k,C_k)$, recovered from the entries
$(B_\mu)_{00}=2M_0$, $(B_\mu)_{k0}=S_k/(\pi k)$, $(B_\mu)_{kk}=2C_{|k|}$; the
quotient has dimension exactly $2N+1$. This is the source-to-matrix
bijection of \cite[Prop.~4.2]{ConnesvanSuijlekom2025}, specialized to the
single-source family used here.
\end{lemma}

\begin{proof}
With the convention $\sgn(0)=0$, for $m\ne n$
$(B_\mu)_{mn}=(\sgn(m)S_{|m|}-\sgn(n)S_{|n|})/(\pi(m-n))$, and
$(B_\mu)_{00}=2M_0$, $(B_\mu)_{kk}=2C_{|k|}$; every entry is linear in the listed
coordinates. The profiles $\omega$, $\sin(2\pi k\omega)$,
$\omega\cos(2\pi k\omega)$ ($1\le k\le N$) are linearly independent (exponential
polynomials), and the displayed entries invert the coordinates, so the quotient
is exactly $2N+1$ dimensional.
\end{proof}

By Lemma~\ref{lem:odd} the profiles carry only odd jets; define the scaled odd
edge jets $J_\ell(f)=\frac{(-1)^\ell}{(2\pi)^{2\ell}}f^{(2\ell+1)}(0)$. A direct
computation gives $J_\ell(f)=a\delta_{\ell,0}+\sum_{k=1}^N(A_k+B_k(2\ell+1))(k^2)^\ell$.

\begin{theorem}[Event-jet determinant]\label{thm:det}
Let $M_N$ be the $d\times d$ matrix, $d=2N+1$, with rows $\ell=0,\dots,2N$ and
columns $e_0,X_1,\dots,X_N,Y_1,\dots,Y_N$ given by
$M_N(\ell,e_0)=\delta_{\ell,0}$, $M_N(\ell,X_k)=(k^2)^\ell$,
$M_N(\ell,Y_k)=(2\ell+1)(k^2)^\ell$. Then
\[
 \det M_N=(-1)^{N(N-1)/2}\,2^N\prod_{k=1}^N k^6
 \prod_{1\le i<j\le N}(j^2-i^2)^4 .
\]
In particular $\dim\mathcal V_N=2N+1$ and the first $2N+1$ scaled odd edge jets
$J_0,\dots,J_{2N}$ are coordinates on $\mathcal V_N$.
\end{theorem}

\begin{proof}
Expand along $e_0$; in the remaining rows set $r=\ell-1$. For node $x_k=k^2$ the
two columns are $x_k^{r+1}$ and $(2r+3)x_k^{r+1}$; replacing the second by
itself minus thrice the first gives $2r x_k^{r+1}$, and factoring $x_k$ and
$2x_k^2$ extracts $2^N\prod_k x_k^3=2^N\prod_k k^6$. The remaining paired-column
matrix is the confluent Vandermonde with each node doubled
\cite{Gautschi1962}, of determinant $\prod_{i<j}(x_j-x_i)^4$. Reordering paired
columns to grouped order costs $N(N-1)/2$ swaps; substituting $x_k=k^2$ gives the
formula.
\end{proof}

\begin{lemma}[Odd edge transport]\label{lem:transport}
For $u_0>0$, $\tau=u-u_0$, $\omega=\tau/(u_0+\tau)$, and integers $r,j\ge1$,
$[\tau^r]\omega^j=(-1)^{r-j}\binom{r-1}{j-1}u_0^{-r}$ for $r\ge j$ and $0$
otherwise. On the odd Taylor-jet subspace (orders $1,3,\dots,4N+1$) the map from
$\omega$-coefficients to $\tau$-coefficients is lower triangular with nonzero
diagonal and determinant $u_0^{-(2N+1)^2}$, and likewise for the scaled jets
$J_0,\dots,J_{2N}$.
\end{lemma}

\begin{proof}
$\omega^j=\tau^ju_0^{-j}(1+\tau/u_0)^{-j}$, and the binomial expansion gives the
coefficient; the diagonal ($r=j$) is $u_0^{-j}$, and
$\sum_{j\ \mathrm{odd}\le 4N+1}j=(2N+1)^2$.
\end{proof}

\begin{proposition}[Observed event stream]\label{prop:stream}
At a detected edge $u_0=\log q$, the odd matrix jumps $\Delta Q_N^{(1)},
\Delta Q_N^{(3)},\dots,\Delta Q_N^{(4N+1)}$ determine, for every entry
$(m,n)$, the scaled stream $J_0(f_{mn}),\dots,J_{2N}(f_{mn})$ with
$f_{mn}=(A_N)_{mn}$; the entries $(0,0)$, $(k,0)$, $(k,k)$ carry the normalized
basis $2b_0,2b_k^s,2b_k^c$ of the finite quotient.
\end{proposition}

\begin{proof}
By Theorem~\ref{thm:rigidity} each odd jump is $(\Lam(q)/\sqrt q)$ times the
corresponding $\tau$-derivative of $f_{mn}(\tau/(u_0+\tau))$; after dividing by
the scalar (known once $\Lam(q)$ is read from the first jump) and inverting the
triangular transport of Lemma~\ref{lem:transport}, one recovers the odd
$\omega$-coefficients of $f_{mn}$, and hence $J_\ell(f_{mn})$ by the fixed
diagonal scaling. The entry profiles of $(B_\mu)_{00}$, $(B_\mu)_{k0}$,
$(B_\mu)_{kk}$ are $2b_0,2b_k^s,2b_k^c$ by Lemma~\ref{lem:quotient}.
\end{proof}

\begin{theorem}[Finite dictionary: window, blind line, recurrence]\label{thm:dictionary}
Let $m_\ell=a\delta_{\ell,0}+\sum_{k=1}^N(A_k+B_k(2\ell+1))(k^2)^\ell$. Then:
\begin{enumerate}
\item[(i)] the first $2N+1$ values $m_0,\dots,m_{2N}$ recover
$(a,A_1,\dots,A_N,B_1,\dots,B_N)$ by exact inversion of $M_N$;
\item[(ii)] the first $2N$ values leave a one-dimensional blind line, on which
$m_{2N}$ is nonzero (the window $2N{+}1$ is sharp);
\item[(iii)] the whole stream is annihilated by
$P_N(S)=S\prod_{k=1}^N(S-k^2)^2$, $(Sm)_\ell=m_{\ell+1}$, and this order is sharp
on the full quotient.
\end{enumerate}
\end{theorem}

\begin{proof}
(i) is Theorem~\ref{thm:det}. For (ii), the $d$ rows of $M_N$ are independent, so
the first $d-1$ have a one-dimensional nullspace; a vector in it with $m_{2N}=0$
would lie in $\ker M_N$, impossible. For (iii), $S$ annihilates $\delta_{\ell,0}$
and, for $x=k^2$, $(S-x)x^\ell=0$ and $(S-x)^2((2\ell+1)x^\ell)=0$; a shorter
recurrence would make the first $2N+1$ rows of $M_N$ dependent, contradicting
Theorem~\ref{thm:det}.
\end{proof}

Thus the source-to-jet map is a lossless finite dictionary with prescribed nodes
$0,1^2,\dots,N^2$ and sharp window $2N+1$ (Figure~\ref{fig:reconstruction}). This
is the finite Prony structure \cite{Prony1795,Gautschi1962,BatenkovYomdin2013} of
the specific matrix $A_N$; by rigidity it recovers, at a prime event, the fixed
profile and the scalar $\Lam(q)$, not further arithmetic.

\begin{figure}[t]
\centering
\includegraphics[width=\textwidth]{fig_reconstruction.pdf}
\caption{The finite source-to-jet dictionary at level $N=4$. Left: the
annihilator nodes of $P_N(S)=S\prod_{k=1}^N(S-k^2)^2$, one simple node at $0$
and $N$ double nodes at $1^2,\dots,N^2$, totalling $2N{+}1$ modes. Right: the
number of source coordinates pinned down by the first $k$ odd jets saturates at
$2N{+}1=9$; one jet short leaves the sharp one-dimensional blind line.}
\label{fig:reconstruction}
\end{figure}

\section{A finite vanishing-moment ceiling at the prime edge}\label{sec:uncertainty}

How blind can a band-limited probe be to an entering prime? For a real even
vector $x=(x_m)_{m\in I_N}$, let $T_x(t)=\sum_m x_m e^{2\pi imt}$, let
$M_j(x)=\sum_m m^jx_m$ be its moments (odd moments vanish by evenness), and let
$e(x)=\min\{j\ge0:M_j(x)\ne0\}$ be its first nonvanishing (even) moment. The
Rayleigh response of $x$ to the entering atom, $x^tA_N(\omega)x=K_x(\omega)$ with
$K_x(\omega)=2\int_0^\omega T_x(t)T_x(\omega-t)\dd t$, has vanishing order
$2e(x)+1$ in $\omega$, hence, via Lemma~\ref{lem:transport}, in the cutoff
variable $\tau$; we call $2e(x)+1$ the \emph{visibility order} of $x$. The
bound below is a vanishing-moment (sum-rule) ceiling; the phrase
\emph{prime-edge uncertainty principle} is kept only as a gloss for its
band-limited reading.

\begin{theorem}[Finite vanishing-moment ceiling at the prime edge]\label{thm:uncertainty}
For every nonzero real even $x$ on $I_N$ the visibility order satisfies
$2e(x)+1\le 4N+1$. Equality holds if and only if, up to a nonzero scalar, $x$ is
the centered finite-difference stencil
\[
        x_m=(-1)^{N-m}\binom{2N}{N+m},\qquad -N\le m\le N,
\]
which is therefore the unique maximally blind band-limited normal, blind to order
$4N+1$. Moreover $K_x(\omega)=\dfrac{2(2\pi)^{2e}M_e(x)^2}{(2e+1)!}\,
\omega^{2e+1}+O(\omega^{2e+3})$, $e=e(x)$.
\end{theorem}

\begin{proof}
The Volterra expansion gives
$[\omega^{k+1}]K_x(\omega)=\frac{2(2\pi i)^k}{(k+1)!}\sum_{a=0}^kM_a(x)M_{k-a}(x)$;
with odd moments zero and $e=e(x)$ the first nonzero even moment, the leading
term is at $k=2e$, giving order $2e+1$ and the stated coefficient: since all odd
moments vanish, $e$ is even, so $i^{2e}=(-1)^e=1$ and the coefficient is real and
positive. The even moments $M_0,M_2,\dots,M_{2N}$ of an even vector are the image
of $(x_0,\dots,x_N)$ under a Vandermonde matrix in the squared nodes
$\{0^2,1^2,\dots,N^2\}=\{0,1,4,\dots,N^2\}$, which is invertible; hence they
cannot all vanish for $x\ne0$, so $e(x)\le 2N$. Requiring $M_0=\cdots=M_{2N-2}=0$ imposes $N$ independent
conditions on the $N+1$ even degrees of freedom, leaving a one-dimensional
solution; the centered stencil satisfies them with $M_{2N}=(2N)!\ne0$, and is
therefore the unique extremizer up to scale. This is the classical
vanishing-moment (sum-rule) bound \cite{Daubechies1992}, here realized as the
prime-edge visibility ceiling of the finite \cvs\ path.
\end{proof}

Figure~\ref{fig:uncertainty} shows the extremal stencil and the visibility
ceiling. The ceiling is sharp and the extremizer unique: a level-$N$
band-limited probe cannot be blind to an entering prime beyond order $4N+1$, and
exactly one probe (the $2N$-th central difference) attains that bound.

\begin{figure}[t]
\centering
\includegraphics[width=\textwidth]{fig_uncertainty.pdf}
\caption{The finite vanishing-moment ceiling at the prime edge, shown at $N=4$. Left: the unique
maximally blind normal, the centered finite-difference stencil
$x_m=(-1)^{N-m}\binom{2N}{N+m}$. Right: the visibility order $2e(x)+1$ of a
family of even normals climbs by $4$ with each additional vanishing moment and
is capped at $4N+1=17$, attained only by the central-difference stencil.}
\label{fig:uncertainty}
\end{figure}

The integer ceiling has a metric refinement. On a band-limited window the leading
term of Theorem~\ref{thm:uncertainty} controls the response, and converts the
visibility \emph{order} into a visibility \emph{scale}.

\begin{corollary}[Edge resolution scale]\label{cor:resolution}
Fix a detected edge $u_0=\log q$ and a nonzero real even $x$ with first
nonvanishing even moment $e=e(x)$. Then there is $\omega_*(x)>0$ on which $K_x$ is
strictly positive and strictly increasing, and, transporting to the cutoff
variable $\tau=u-u_0$ by $\omega=\tau/(u_0+\tau)$, the advance past the edge at
which $|K_x|$ first reaches a threshold $\varepsilon$ satisfies, as
$\varepsilon\to0$,
\[
        \tau_\varepsilon(x)=\log q\,
        \Bigl(\tfrac{(2e+1)!\,\varepsilon}{2(2\pi)^{2e}M_e(x)^2}\Bigr)^{1/(2e+1)}
        (1+o(1)).
\]
The first-passage exponent $1/(2e(x)+1)$ decreases with the visibility order, so
for small $\varepsilon$ the maximally blind probe $e=2N$ has the coarsest edge
resolution and the generic probe $e=0$ the finest: the integer ceiling of
Theorem~\ref{thm:uncertainty} is the order of this resolution scale.
\end{corollary}

\begin{proof}
By Theorem~\ref{thm:uncertainty},
$K_x(\omega)=\tfrac{2(2\pi)^{2e}M_e^2}{(2e+1)!}\omega^{2e+1}(1+O(\omega^2))$ with a
positive leading coefficient. The remainder is an analytic multiple of
$\omega^{2e+3}$, so it is dominated by the leading term on a nonempty window
$0<\omega\le\omega_*(x)$ (such a window exists because the remainder is
$O(\omega^{2e+3})$ with the leading term of exact order $\omega^{2e+1}$; the
window scale is set by the ratio of successive even moments of $x$);
there the leading term controls sign, monotonicity, and, as $\varepsilon\to0$,
the first-passage scale. Inverting $K_x(\tau)=\varepsilon$ with
$\omega=\tau/\log q+O(\tau^2)$ gives $\tau_\varepsilon$. The window $\omega_*(x)$
depends on $x$ through those moment ratios and is not uniform over the unit
sphere.
\end{proof}

\begin{remark}
The integer $2e(x)+1$ is upper semicontinuous and generically equal to $1$: every
neighborhood of the extremal stencil contains probes of visibility order $1$, so
the ceiling is attained only on a measure-zero set. The resolution scale
$\tau_\varepsilon$ is the robust, metric form of the same statement.
\end{remark}

\section{The coincidence-resolvent generating identity}\label{sec:deflation}

The first-derivative jump is a rank-one negative shock in the fixed universal
direction $\one_N$. Contracting the resolvent against that direction generates,
in one rational function, the jump of every spectral flow: the log-determinant,
every trace moment, the eigenvalue velocities, and the elementary barriers.

\begin{theorem}[Coincidence-resolvent generating identity]\label{thm:generating}
Let $u_0=\log q$ be a prime-power event, $Q_0=Q_N(u_0)$, and
$a_q=2\Lam(q)/(\sqrt q\,\log q)>0$. Define the \emph{coincidence resolvent}
\[
        G_N(z)=\bigl\langle\one_N,(Q_0-z)^{-1}\one_N\bigr\rangle
        =\sum_{j=1}^d\frac{|\langle\phi_j,\one_N\rangle|^2}{\lambda_j-z},
        \qquad z\notin\operatorname{spec}(Q_0),
\]
a rational function whose poles lie among the eigenvalues of $Q_0$: grouping the
sum by spectral projections $P_\lambda$,
$G_N(z)=\sum_{\lambda\in\operatorname{spec}(Q_0)}\|P_\lambda\one_N\|^2/(\lambda-z)$
has a simple pole with residue $-\|P_\lambda\one_N\|^2$ at exactly those
eigenvalues $\lambda$ with $P_\lambda\one_N\ne0$; an eigenvalue whose eigenspace
is orthogonal to $\one_N$ contributes no pole (the corresponding point is
removable), and a repeated eigenvalue contributes one grouped residue, not one
pole per basis eigenvector. Then, at $u_0$, for $z\notin
\operatorname{spec}(Q_0)$ (only the derivative jump enters, not the branch of
$\log\det$),
\[
 \Delta\frac{d}{du}\log\det(Q_N(u)-z)=-a_q\,G_N(z),
 \qquad
 \Delta\frac{d}{du}\operatorname{tr}Q_N(u)^p=-p\,a_q\,
 \bigl\langle\one_N,Q_0^{\,p-1}\one_N\bigr\rangle\ \ (p\ge1).
\]
In particular $\Delta\frac{d}{du}\operatorname{tr}Q_N(u_0)=-a_q\,d$, so the scalar
trace path $\operatorname{tr}Q_N$ has singular curvature
$(D_u^2\operatorname{tr}Q_N)_{\rm sing}=-d\sum_q a_q\delta_{\log q}$: its
singular second
derivative is exactly $-(2N+1)$ times the scalar von Mangoldt measure
$\mu_{mn}$ of Corollary~\ref{cor:measure}, hence carries the prime side of the
explicit formula under the same conversion.
\end{theorem}

\begin{proof}
$Q_N$ is continuous at $u_0$, so $(Q_N(u)-z)^{-1}$ and $Q_N(u)^{p-1}$ are
two-sided continuous there; the only discontinuity is the velocity jump
$\Delta Q_N'(u_0)=-a_q\one_N\one_N^t$ (Theorem~\ref{thm:jump}). By Jacobi's
formula $\frac{d}{du}\log\det(Q_N-z)=\operatorname{tr}((Q_N-z)^{-1}Q_N')$, whose
jump is $\operatorname{tr}((Q_0-z)^{-1}\Delta Q_N')=-a_q\one_N^t(Q_0-z)^{-1}\one_N$;
the spectral decomposition gives $G_N$. Likewise
$\frac{d}{du}\operatorname{tr}Q_N^p=p\operatorname{tr}(Q_N^{p-1}Q_N')$ has jump
$-p\,a_q\,\one_N^tQ_0^{\,p-1}\one_N$; for $p=1$ this is $-a_q\|\one_N\|^2=-a_qd$.
\end{proof}

The residues of $G_N$ are $-|\langle\phi_j,\one_N\rangle|^2$; their magnitudes are
the couplings to the universal direction, and these magnitudes, scaled by
$-a_q$, are exactly
the eigenvalue-velocity jumps of Corollary~\ref{cor:deflation}.

\begin{remark}[The coincidence spectral measure]
Writing $\rho=\sum_j|\langle\phi_j,\one_N\rangle|^2\,\delta_{\lambda_j}$ for the
finite positive measure carried by the couplings, $G_N(z)=\int(\lambda-z)^{-1}\dd\rho$
is its Cauchy transform and every trace-moment jump is a moment of $\rho$:
$\Delta\frac{d}{du}\operatorname{tr}Q_N^p=-p\,a_q\int\lambda^{\,p-1}\dd\rho$, with
$\rho(\R)=\|\one_N\|^2=2N+1$ and $\int\lambda\,\dd\rho=\langle\one_N,Q_0\one_N\rangle$.
Thus a single finite measure governs all the spectral-flow jumps at the event.
\end{remark}

\begin{corollary}[Velocity deflation and interlacing]\label{cor:deflation}
The velocity spectrum deflates by rank one: $Q_N'(u_0{+}0)=Q_N'(u_0{-}0)
-a_q\one_N\one_N^t$, so (Weyl/Cauchy) the right velocity spectrum interlaces the
left,
\[
 \lambda_1\!\bigl(Q_N'(u_0{+}0)\bigr)\le\lambda_1\!\bigl(Q_N'(u_0{-}0)\bigr)
 \le\lambda_2\!\bigl(Q_N'(u_0{+}0)\bigr)\le\cdots
 \le\lambda_d\!\bigl(Q_N'(u_0{-}0)\bigr),
\]
and, for simple $Q_0$, each eigenvalue velocity decelerates by
$\Delta\lambda_j'(u_0)=-a_q|\langle\phi_j,\one_N\rangle|^2\le0$ (the residues of
$G_N$), with total $\sum_j\Delta\lambda_j'=-a_qd$, strict off
$\one_N^{\perp}$.
\end{corollary}

\begin{proof}
Immediate from $\Delta Q_N'(u_0)=-a_q\one_N\one_N^t$ and, for the velocities,
first-order perturbation $\Delta\lambda_j'=\langle\phi_j,\Delta Q_N'\phi_j\rangle$.
\end{proof}

\begin{remark}[Elementary spectral barriers in a positive cell; conditional]\label{thm:barriers}
Suppose $Q_0=Q_N(u_0)$ is positive definite. (This is a hypothesis, not a
verified property of the actual \cvs\ path at the program's canonical
parameters, where the assembled form has near-zero eigenvalues
\cite{Groskin2026Paper1}; whether any
actual prime-power event sits in a positive cell is not established here, and
the statement is conditional on it, which is why it is recorded as a remark
rather than a theorem.) Let $e_r(Q)$ be the $r$th elementary
symmetric function of the eigenvalues of $Q$. Then for every $1\le r\le d$,
\[
        \Delta\frac{d}{du}e_r(Q_N(u))\Big|_{u=u_0}
        =-a_q\,\one_N^tT_{r-1}(Q_0)\one_N<0,\qquad
        T_{r-1}(Q_0)=\sum_{j=0}^{r-1}(-1)^je_{r-1-j}(Q_0)Q_0^j .
\]
In particular the determinant barrier ($r=d$) and its logarithm strictly
decelerate at every prime event in such a cell.
\end{remark}

\begin{proof}[Proof of the barrier inequalities]
For a principal $r$-set $S$, the matrix determinant lemma gives
$\det(Q_{0,S}-ta_q\one_S\one_S^t)=\det Q_{0,S}-ta_q\,\one_S^t\adj(Q_{0,S})\one_S$;
$Q_{0,S}$ is positive definite so $\adj(Q_{0,S})$ is positive definite and each
summand is positive; summing over $|S|=r$ proves strict negativity. The
Newton-transformation identity $De_r(Q)[H]=\operatorname{tr}(T_{r-1}(Q)H)$ with
$H=-a_q\one_N\one_N^t$ gives the displayed value; in an eigenbasis of $Q_0$,
$T_{r-1}(Q_0)$ is diagonal with positive entries, so the quadratic form on
$\one_N$ is positive.
\end{proof}

\subsection*{The prime event as a rank-one Weyl-function increment}

The rank-one velocity deflation has a clean resolvent reading. A related
scalar-resolvent reduction appears in Andrade's analysis of the pole term of
the continuum localized form \cite{Silva20682834}, where a
rank-one-per-parity splitting reduces a bottom-eigenvalue question to a scalar
resolvent criterion with a Krein analysis; the object here is different, namely
the prime-event velocity jump of the fixed-$N$ path rather than the pole term.
Let
$V_\pm=Q_N'(u_0\pm0)$ be the
right and left spectral velocities at $u_0=\log q$, so $V_+=V_--a_q\one_N\one_N^t$
by Theorem~\ref{thm:jump}, and define the \emph{coincidence-Weyl function} of the
velocity in the universal direction,
\[
        W_\pm(z)=\bigl\langle\one_N,(V_\pm-z)^{-1}\one_N\bigr\rangle,
        \qquad z\notin\operatorname{spec}(V_\pm).
\]
$W_\pm$ is a rational Nevanlinna function; when $\one_N$ is cyclic for $V_\pm$
it is $(2N+1)$ times the Weyl function of the finite Jacobi matrix obtained by
tridiagonalizing
$V_\pm$ from the normalized vector $\one_N/\sqrt{2N+1}$ (the identities below
need no such assumption).

\begin{theorem}[Prime event as a rank-one Weyl increment]\label{thm:krein}
For every $z$ off the spectra of $V_\pm$, $1-a_qW_-(z)\ne0$ and
\[
        W_+(z)=\frac{W_-(z)}{1-a_qW_-(z)},
        \qquad a_q=\frac{2\Lam(q)}{\sqrt q\,\log q}>0 .
\]
Equivalently, as an identity of meromorphic functions,
\[
        \frac{1}{W_+(z)}=\frac{1}{W_-(z)}-a_q,
\]
valid pointwise exactly where $W_+(z)W_-(z)\ne0$; the constant decrement is
\emph{independent of $z$}. (The pointwise restriction is not vacuous: a
nonconstant real rational Nevanlinna function has real zeros interlacing its
poles, so $W_\pm$ generically vanish at points off both spectra whenever
$\one_N$ has
nonzero projection onto at least two eigenspaces, the exceptional case being an
interlacing zero that coincides with an eigenvalue whose eigenspace is
orthogonal to $\one_N$, and at such points off the spectra both
reciprocals are undefined.) Equivalently $V_+$ is the Schur complement, with
respect to the scalar block, of the bordered matrix
$\left(\begin{smallmatrix}a_q^{-1}&\one_N^t\\ \one_N&V_-\end{smallmatrix}\right)$,
i.e.\ a single symmetric elimination step with pivot $a_q^{-1}=\sqrt q\,\log q/
(2\Lam(q))$. Consequently the von Mangoldt weight is read off from the
reciprocal increment,
\[
        \Lam(q)=\frac{\sqrt q\,\log q}{2}\Bigl(\tfrac{1}{W_-(z)}-\tfrac{1}{W_+(z)}\Bigr)
        \qquad\text{wherever }W_+(z)W_-(z)\ne0 .
\]
\end{theorem}

\begin{proof}
For $z$ off both spectra,
$\det(V_+-z)=\det(V_--z)\,(1-a_qW_-(z))$ by the matrix determinant lemma, so
$1-a_qW_-(z)\ne0$ there. The Sherman--Morrison identity applied to
$V_+=V_--a_q\one_N\one_N^t$ with
$M=V_--z$ gives
$W_+=\one_N^t(M-a_q\one_N\one_N^t)^{-1}\one_N=W_-/(1-a_qW_-)$. Dividing by
$W_+W_-$ where both are nonzero yields the reciprocal form, which then extends
to the stated meromorphic identity. The Schur statement is the identity
$V_--\one_N\,a_q\,\one_N^t=V_+$ for the displayed border. The zero/pole
interlacing for a real rational Nevanlinna function is classical
\cite{KacKrein1974R}
($W_-'>0$ between consecutive real poles).
\end{proof}

\begin{remark}[Krein-string analogy, not established here]
The constant decrement of $1/W$ is \emph{formally} the boundary point-mass move
in Krein's string--to--continued-fraction correspondence
\cite{KacKrein1974string}, and it is tempting to
read $a_q$ as a per-event boundary mass in the same de Branges--Krein
canonical-system setting as Suzuki's continuous screw-function pencil
\cite{Suzuki2026}, with the reduced parametrization ($\omega_q'(u_0)=1/\log q$)
matching the prime data his continuous screw function carries. Making that
reading literal, however, requires Stieltjes-class hypotheses (a positive
string with spectral measure supported on a half-line
\cite{KacKrein1974string}) that the indefinite
velocities $V_\pm$ are not proved to satisfy here; a general rational Nevanlinna
function is the Weyl function of a finite Jacobi matrix or canonical system
\cite{GesztesySimon1997,KacKrein1974R},
which is a strictly larger class than the Krein strings. We therefore record
the string reading as an analogy, not a theorem. In any case only the
\emph{per-event} reciprocal increment equals $a_q$; the interior structure of
$V_-$ encodes its accumulated spectral history, not the individual primes, and
nothing here bears on positivity of the finite form or on the zeros.
\end{remark}

\section{Universality, the explicit formula, and scope}\label{sec:scope}

\paragraph{Universality across the Selberg class.}
The transfer below is conditional on the construction of the corresponding
finite path: it applies to any finite path assembled with the same prime-block
template as the \cvs\ path, with the $\zeta$ source weights replaced by those of
$L(s,\pi)$; such a path is exhibited here only for $\zeta$. For Dirichlet
twists the prime block is written down entrywise below, but no full twisted
path (in particular no twisted archimedean or pole block) is constructed, and
the construction for a general Selberg-class $L$-function
is not carried out in this paper. Granting the path, the prime block is
determined only by the source amplitudes and the
test-function-at-$\log q$ pairing: the single-source matrix $A_N(1-\log q/u)$ is
the same divided-difference matrix regardless of the $L$-function, since the
gamma factor enters only the archimedean block $Q_{\mathrm{arch}}$, not the prime
term. For an $L$-function $L(s,\pi)$ with prime side
$\sum_q c_\pi(q)q^{-1/2}\phi(\log q)$
(Dirichlet $c_\chi(q)=\Lam(q)\chi(q)$, or automorphic $c_\pi(q)=\Lam(q)a_\pi(q)$),
such a finite path has first-derivative jump
$-\,2c_\pi(q)(\log q)^{-1}q^{-1/2}\one_N\one_N^t$, and the whole event
\emph{algebra} transfers verbatim: the rank-one direction $\one_N\one_N^t$, the
source-to-jet dictionary on the squared nodes $\{0,1^2,\dots,N^2\}$, and the
vanishing-moment ceiling depend only on $A_N$ and hold for arbitrary (in particular
complex) weights. The \emph{sign-definite} consequences are different: negative
semidefiniteness of the event measure (Corollary~\ref{cor:measure}), the velocity
interlacing and deceleration (Corollary~\ref{cor:deflation}), and the strict
barrier decrease (Remark~\ref{thm:barriers}) all use $a_q>0$, hence a
\emph{positive} weight $c_\pi(q)>0$. This holds at every prime power for $\zeta$
(where $c(q)=\Lam(q)>0$), so those results transfer verbatim there, the barrier
statement remaining conditional on its positive-definiteness hypothesis
(Remark~\ref{thm:barriers}). For a real
primitive character the atom is Hermitian, but $c(q)=\Lam(q)\chi(q)$ is negative
wherever $\chi(q)=-1$ (for instance $\chi_4(3)=-1$), and there the three
sign-definite statements reverse; they hold only at prime powers with
$\chi(q)>0$. For complex $c_\pi(q)$ the atom is not even Hermitian. For Dirichlet twists the family of first jumps over
$\chi\bmod m$ separates the edge by residue class through character orthogonality,
a finite rendering of the prime side of the twisted explicit formula
\cite{IwaniecKowalski2004}.

\begin{corollary}[Residue-class readout across the Dirichlet family]\label{cor:dirichlet}
Fix a prime-power edge $u_0=\log q$ with $(q,m)=1$. Assume, for each Dirichlet
character $\chi\bmod m$, a finite path with the same prime-block template,
carrying the twisted weights $c_\chi(q')=\Lam(q')\chi(q')$, whose non-prime
blocks are real analytic at $u_0$; such paths are granted, not constructed,
here. Then the first-derivative jump of the $\chi$-path is
$J^{(\chi)}=j^{(\chi)}\one_N\one_N^t$ with
$j^{(\chi)}=-2\Lam(q)\chi(q)(\log q)^{-1}q^{-1/2}$ (by the argument of
Theorem~\ref{thm:jump}), and for every residue $a$ with
$(a,m)=1$,
\[
 -\frac{\sqrt q\,\log q}{2}\cdot\frac{1}{\phi(m)}
 \sum_{\chi\bmod m}\overline{\chi(a)}\,j^{(\chi)}
 =\Lam(q)\,\mathbf 1[\,q\equiv a\ (\mathrm{mod}\ m)\,].
\]
The already-localized weight $\Lam(q)$ is thus resolved by residue class from the
joint first jumps of the Dirichlet family; this reads $\Lam(q)$ in an arithmetic
progression at a known edge, and is not a prediction of $q$.
\end{corollary}

\begin{proof}
The jump formula follows exactly as in Theorem~\ref{thm:jump}: the analytic
non-prime blocks cancel in the right-minus-left derivative, and the entering
atom carries the twisted weight. Then
substitute $j^{(\chi)}$ and apply the character orthogonality
$\frac{1}{\phi(m)}\sum_{\chi\bmod m}\overline{\chi(a)}\chi(q)
=\mathbf 1[q\equiv a\ (\mathrm{mod}\ m)]$, valid for $(q,m)=1$.
\end{proof}

\paragraph{The prime side of the explicit formula.}
Corollary~\ref{cor:measure} places the primes as the singular second derivative
of $Q_N$: the finite path carries the \emph{prime side} of the Weil--Guinand
explicit formula \cite{Weil1952,IwaniecKowalski2004} exactly and without an
archimedean cutoff, as a matrix-valued measure. This is by construction (the
prime sum is a defining summand of $Q_N$), and we prove nothing here about the
zero side: the identity relating the prime side to the zeros is the classical
explicit formula, supplied externally by the archimedean and pole blocks of
Paper~2 \cite{Groskin2026FiniteDictionary} and the eigenvalue-to-zero
approximation of Paper~1 \cite{Groskin2026Paper1} (with a quantitative
convergence law for the first zero in the CCM setting in
\cite{Andrews2026}). The object is thus
complementary to Suzuki's continuous screw-function picture \cite{Suzuki2026},
to the CCM spectral triple \cite{ConnesConsaniMoscovici2025}, and to Liu's
determinant encoding of local Euler factors \cite{Liu2026}, each a differently
mechanized realization of the same arithmetic.

\paragraph{Scope.}
The claims are finite and structural. The paper proves no positivity of the
finite form, no Riemann Hypothesis, no prime-counting or next-prime bound, and
no factoring statement; by Theorem~\ref{thm:rigidity} the local event carries
exactly $\Lam(q)$, and the reconstruction, vanishing-moment, and deflation results are
finite harmonic analysis of the fixed path, not arithmetic prediction. The event
jets are right-minus-left jumps taken after the common analytic background has
cancelled.

\paragraph{Verification.}
The proofs are algebraic; the companion scripts are separate reproducibility
checks for
signs, indexing, determinant factors, and finite-field behavior, not proof
substitutes. The ancillary verification package records (scripts and JSON
artifacts listed in its manifest, each script writing the JSON artifact of the
same stem, e.g.\ \texttt{rank\_one\_jump\_audit.json}), over exact integer or
modular arithmetic:
\begin{center}
\footnotesize
\setlength{\tabcolsep}{3.5pt}
\begin{tabular}{lll}
rank-one jump, end-to-end & $q\in\{3,4,5,7,8,9,25\}$, $N\le6$
 & \texttt{check\_rank\_one\_jump}\\
event-jet and transport rank & $N\le200$, four prime fields
 & \texttt{check\_event\_jet\_largeN}\\
event-jet and transport determinants & $N\le200$, four prime fields
 & \texttt{check\_event\_jet\_determinant}\\
recurrence & $N\le1000$, four prime fields
 & \texttt{check\_event\_jet\_recurrence}\\
event dictionary (window, blind line) & $N\le200$
 & \texttt{check\_event\_prony\_reconstruction}\\
source quotient and transport & $N\le200$
 & \texttt{check\_source\_quotient\_and\_transport}\\
vanishing-moment ceiling (extremizer) & $N\le10$
 & \texttt{check\_uncertainty\_ceiling}\\
spectral-barrier deceleration & 50 positive definite cases
 & \texttt{check\_spectral\_barrier\_jump}\\
universal jet $B_{r,N}$ closed form & exact symbolic, orders $r\le5$
 & \texttt{check\_universal\_jet}\\
coincidence-averaged readout & $(2N+1)^2$ variance (IID model)
 & \texttt{check\_coincidence\_readout}\\
Dirichlet residue-class readout & cyclic moduli $m\le13$
 & \texttt{check\_dirichlet\_readout}
\end{tabular}
\end{center}
As a floating-point check on the actual assembled operator at the program's
canonical scale, at $N=200$ (dimension $401$) every prime power $q\le100$ was
confirmed to have a rank-one first-derivative jump equal to
$-2\Lam(q)/(\sqrt q\,\log q)\one_N\one_N^t$: over all $35$ events the second
singular value of the jump is at most $4\times10^{-7}$ of the first, the jump
matches $-a_q\one_N\one_N^t$ to relative error $10^{-6}$, and $\Lam(q)$ is
recovered from a single entry to $5\times10^{-8}$; these four figures are
produced by \texttt{check\_canonical\_scale.py} and recorded in its artifact
\texttt{canonical\_scale\_audit.json} in the ancillary archive. The archimedean and pole
blocks (Lemma~\ref{lem:background}) are analytic and cancel in the jump, so the
prime part carries the entire singular event.

\paragraph{Statement on use of AI tools.}
Language-model tools were used as a drafting and verification-tooling assistant
under the author's direction; the author takes sole responsibility for all
mathematics, computations, and text.

\paragraph{Data and code availability.}
The verification package (the separate reproducibility check scripts, their exact JSON
artifacts, and the figure generator) accompanies this paper as ancillary files;
a checksummed archive is deposited on Zenodo, concept DOI
\href{https://doi.org/10.5281/zenodo.21242028}{10.5281/zenodo.21242028} (which
resolves to the latest version). This work is licensed under CC-BY-4.0.

{\footnotesize\raggedright
\bibliographystyle{plainurl}
\bibliography{main}
}

\end{document}
